\documentclass[12pt,a4paper]{article}
\usepackage{fontspec}
\setmainfont{Times New Roman}
\usepackage[top=2cm,bottom=2cm,left=2.5cm,right=2cm]{geometry}
\usepackage{enumitem}
\usepackage{tabularx}
\usepackage{booktabs}
\usepackage{array}
\newcolumntype{Y}{>{\centering\arraybackslash}X}
\pagestyle{plain}
\linespread{1.15}

\begin{document}
% --- TÀI LIỆU VẬT CHỨNG HIỆN TRƯỜNG ---
\begin{center}
    \fontsize{11pt}{13pt}\selectfont
    \textbf{HỒ SƠ VẬT CHỨNG — CHUYÊN ÁN MẠNG SỐ 14 ĐƯỜNG BỜ SÔNG}\\[-0.2em]
    \rule{6cm}{0.6pt}\\[0.5cm]
    \textbf{\Large  VẬT CHỨNG: TRÍCH ĐOẠN BÁO CŨ NĂM 1996 THU THẬP TẠI HIỆN TRƯỜNG}\\[0.15cm]
    \textit{}
\end{center}
\vspace{0.4cm}

\subsection*{VỤ ÁN: CHUYÊN ÁN \#000 — TRỐN TÌM}


> \textbf{Ký hiệu vật chứng hiện trường:} Cọc tiêu số \texttt{05} (Thu giữ tại nền nhà phòng khách, quanh chân nạn nhân Nguyễn Văn Khang).

> \textbf{Hình thức vật chứng:} Các mảnh báo in đen trắng ố vàng bị xé vụn (6–8 mảnh).

> \textbf{Hướng dẫn đạo cụ thực tế:} Bản in tài liệu dưới đây được in ra giấy A4/A5, có thể ngâm nước chè loãng sấy khô để tạo màu ố vàng thập niên 1990, sau đó xé thủ công thành 6–8 mảnh mép răng cưa để người chơi trực tiếp lắp ghép, khám phá manh mối quá khứ.




\begin{quote}
\fontsize{9pt}{11pt}\selectfont\ttfamily
+----------------------------------------------------------------------------------------+\\
| BÁO THỦ ĐÔ THỜI ĐẠI               TIẾNG NÓI CỦA NHÂN DÂN THỦ ĐÔ             SỐ 182/1996 |\\
| Cơ quan của Thành ủy Hà Nội       Phát hành Thứ Sáu, ngày 26/07/1996        Giá: 800đ  |\\
+----------------------------------------------------------------------------------------+\\
\end{quote}



\subsection*{CHUYÊN MỤC: AN NINH \& ĐỜI SỐNG DÂN SINH}



\subsection*{\textbf{BI KỊCH TỪ TRÒ CHƠI TRỐN TÌM: MỘT CHÁU BÉ TỬ VONG TRONG TỦ GỖ}}


\textbf{HÀ NỘI (PV) — Chiều ngày 24 tháng 07 năm 1996, một vụ tai nạn thương tâm bắt nguồn từ trò chơi con trẻ đã cướp đi sinh mạng của một cháu bé 7 tuổi tại khu dân cư ven Bờ Sông, Phường Phân khu Cảng, TP. Hà Nội.}


Theo thông tin ban đầu từ Công an cơ sở, vào khoảng 15 giờ 30 phút chiều ngày 24/07/1996, một nhóm trẻ em xóm Bờ Sông rủ nhau chơi trò trốn tìm tại khu vực nhà kho và gian nhà hoang cuối ngõ. Trong lúc chơi đùa, cháu \textbf{N.G.H} (sinh năm 1989, 7 tuổi, trú tại ngõ Bờ Sông) đã chui vào ẩn nấp bên trong một chiếc tủ gỗ đựng chăn màn cũ kích thước lớn bỏ không. Chiếc tủ này có kết cấu cánh cửa gỗ dày, có then cài bằng sắt lắp ở mặt bên ngoài.


Sau khi cháu N.G.H chui vào trong, một bạn cùng chơi trong nhóm là cháu \textbf{N.V.K} (8 tuổi, trú cùng ngõ) do không để ý và mải chạy trốn đã đẩy sập cánh cửa tủ rồi vô tình gạt chốt then sắt bên ngoài, khóa kín chiếc tủ lại mà không nhận ra có bạn đang trốn bên trong.


Đau xót hơn, cháu N.G.H không may bị tật câm bẩm sinh từ nhỏ, không thể cất tiếng la hét kêu cứu. Chiếc còi đồng thau nhỏ gắn dây dù màu cam được gia đình đeo trước ngực cháu để phòng khi cần báo hiệu đã bị vướng kẹt vào khe nẹp đáy tủ trong lúc cháu ngồi bó gối, khiến cháu không thể thổi còi báo động.


Sau khi tan cuộc chơi vào chiều muộn, nhóm trẻ tản mát về nhà. Đến chập tối, gia đình và người anh trai của cháu là cháu \textbf{N.T.T} (10 tuổi) không thấy em về ăn cơm nên đã cùng bà con xóm Bờ Sông tỏa đi tìm kiếm khắp nơi. Mãi đến 20 giờ 15 phút cùng ngày, khi người anh trai lần tới khu nhà hoang và bật chốt then cửa tủ gỗ, mọi người mới bàng hoàng phát hiện cháu N.G.H đã rơi vào tình trạng hôn mê sâu, ngưng thở. Dù được trạm y tế địa phương tích cực cấp cứu, cháu bé đã không qua khỏi do ngạt khí cấp tính kéo dài.


Nhận được tin báo, lực lượng chức năng và Công an địa phương đã có mặt bảo vệ hiện trường và tiến hành xác minh. Qua quá trình khám nghiệm và lấy lời khai các nhân chứng cùng nhóm trẻ chơi chiều hôm đó, cơ quan chức năng xác định đây là một vụ tai nạn rủi ro con trẻ vô ý ngoài ý muốn, không có dấu hiệu tội phạm hình sự hay cố sát. Chiếc tủ gỗ sau đó đã được niêm phong và tiêu hủy theo yêu cầu.


Sự việc đau lòng trên là hồi chuông cảnh tỉnh sâu sắc gửi tới các bậc phụ huynh và các hộ gia đình trong việc quản lý, trông nom con trẻ trong những ngày hè, tránh để các em tự do chơi đùa tại những nơi hoang vắng, ẩn chứa hiểm họa khôn lường.


\textbf{(Bài và ảnh: Nhóm PV Thời sự - Bạn đọc)}




\subsection*{ DẤU VẾT THU NHẬN KHI GIÁM ĐỊNH MẢNH BÁO TẠI HIỆN TRƯỜNG:}

\begin{itemize}[leftmargin=1.2cm, itemsep=2pt]
  \item \textbf{Vết dính sinh học \& vật lý:}
  \item Tại góc mép rách của mảnh báo số 3 có dính một vệt máu khô nhỏ (nhóm máu O, trùng với nhóm máu của nạn nhân Nguyễn Văn Khang).
  \item Bề mặt các mảnh giấy báo có bám dính nhiều bụi bột vôi vữa xi măng trắng xám (đặc trưng của thợ xây/thợ nề Nguyễn Thanh Tùng).
  \item \textbf{Mối liên hệ nhân vật:}
  \item Cháu bé tử vong \textbf{N.G.H} (7 tuổi): Nguyễn Gia Huy (em trai của thợ nề Nguyễn Thanh Tùng).
  \item Cháu bé vô tình gài then \textbf{N.V.K} (8 tuổi): Nạn nhân Nguyễn Văn Khang lúc nhỏ.
  \item Người anh trai phát hiện thi thể \textbf{N.T.T} (10 tuổi): Nghi phạm Nguyễn Thanh Tùng (người mang bài báo đến đối chất với Khang đúng ngày giỗ 20 năm 24/07/2016 rồi xé vụn).
\end{itemize}
\end{document}
